% 恶补报告 · 版式导言区（经 pandoc --include-in-header 注入）
%
% 【设计约束】本文件**不引入任何新宏包**。
%   理由：pandoc 默认模板已按 -V 参数配置好 fontspec / xeCJK / geometry / hyperref / setspace；
%   再 `\usepackage` 同名包会触发 option clash 而**编译失败**；
%   而引入 MiKTeX 未装的包会触发首次安装（可能失败或长时间阻塞）。
%   故这里只做「长度与间距」类微调 —— 零新依赖，零冲突面。
%
% 【口径单源】字体 / 页边距 / 字号 / 行距由 cram_render.py 的 LAYOUT 通过 -V 传入，
%   本文件**不得**再重复设置这些量（那会造成两处口径）。

\AtBeginDocument{%
  % 中文段落用「段间距」而不是「首行缩进」来分段：报告里段落短、缩进会显得碎
  \setlength{\parindent}{0pt}%
  \setlength{\parskip}{0.5em}%
}

% 引用块（每条知识点末尾的「> 来源：…」）左右留白收窄 ——
% 默认的 quote 环境会让每条来源锚占掉接近半屏，报告会显得松散。
\makeatletter
\renewenvironment{quote}%
  {\list{}{\rightmargin=1.0em\leftmargin=1.4em}\item\relax}%
  {\endlist}
\makeatother
